\chapter{Methodology}

\section{Overview of the approach}

The methodology follows a complete document-processing pipeline designed to
transform heterogeneous textual resources into a structured corpus that can
support search, filtering, comparison, and thematic analysis. The overall
objective is to make dispersed descriptions of territorial initiatives related
to population ageing more accessible for research and future operational use.

The methodological process evolved during the internship. The initial approach
relied on a multi-step rule-based pipeline: documents were cleaned, segmented,
analysed for document sections, and processed to extract candidate metadata
values. Candidate values were then scored according to explicit rules. This
approach was transparent and useful for inspecting intermediate results, but the
diversity of source formats and writing styles made it difficult to generalise
the extraction rules to every document type.

The final approach therefore complements this rule-based pipeline with
LLM-assisted information extraction. A local language model receives a
constrained prompt and transforms each textual document into a structured JSON
record. The methodology is consequently organised around seven components:
collection, corpus construction, cleaning and quality control, progressive
structuring, constrained LLM extraction, post-extraction taxonomy
normalisation, and evaluation.

\begin{figure}[H]
    \centering
    \includegraphics[width=0.82\textwidth]{Figures/index-insight-pipeline.pdf}
    \caption{Architecture of the Index\_Insight pipeline. Heterogeneous
    documentary sources are collected and normalised before being processed by
    a traceable rule-based baseline and by constrained local LLM extraction.
    Structured records are evaluated against a human reference sample and used
    to generate exploratory analytical outputs.}
    \label{fig:index-insight-pipeline}
\end{figure}

\section{Data collection}

Data were collected from multiple sources using dedicated Python scripts. The
collection process focused on web pages, PDF documents, action sheets, and
project descriptions related to ageing, autonomy, social inclusion, housing,
mobility, health, participation, and support for caregivers. Each retained
resource was stored as a separate \texttt{.txt} file, following the principle
of one documentary record per file.

The corpus contains 1,153 textual documents. Sixty documents were provided
directly by the SIAGE Chair as action sheets, while the remaining documents were
collected from external institutional or associative sources. The collection
scripts use HTTP requests and HTML parsing to retrieve the textual content of
web pages. Additional processing is used when a relevant resource is provided
as a PDF document. Since the websites differ in navigation, HTML structure, and
metadata availability, the collection process relies on source-specific
strategies rather than on a single generic scraper.

The final cleaned corpus contains 1,153 documents and 917,554 word tokens,
corresponding to an average of 795.8 words per document. Document length ranges
from 8 to 3,573 words, illustrating the heterogeneous nature of the collected
material. The corpus is therefore composed primarily of medium-length resources
that usually contain sufficient contextual information to support extraction,
while still requiring robust processing because relevant information may occur
at different locations and be expressed in highly variable forms.

The collection process preserves a link between each record and its source.
This provenance is necessary for quality control, source attribution, correction
of extraction errors, and future updates. The project is intended for
documentary and exploratory purposes; collection should therefore continue to
respect the terms of access of the relevant websites and avoid excessive request
rates.

For every LLM-generated record, the final index also stores explicit provenance
metadata: \texttt{source\_site}, \texttt{source\_url},
\texttt{original\_title}, and \texttt{collection\_date}. These fields are
populated deterministically from the collection indexes and archived source
files; they are not generated by the language model. The source URL points to
the original web page or PDF when it is available, whereas the internal SIAGE
action sheets are identified as such and do not receive an artificial public
URL. The original title preserves the wording supplied by the source and makes
it possible to distinguish the source document from a normalised or extracted
metadata value.

At the time of analysis, all 1,143 available LLM records contained a
standardised source identifier. A source URL was available for 1,083 external
records; the remaining 60 records correspond to SIAGE internal action sheets.
An original title was available for 1,142 records, with one RFVAA record lacking
a title in its source index. These coverage figures make missing provenance
visible rather than replacing it with generated or inferred information.

The historical collection indexes did not consistently record a reliable
collection timestamp. The \texttt{collection\_date} field is therefore left
empty for these legacy records rather than being reconstructed from file-system
dates, which would be misleading. Future collection scripts should record the
timestamp automatically at download time, using an ISO 8601 date. Together,
these provenance fields support attribution, verification against the original
document, correction of metadata, and responsible reuse of externally collected
materials.

\section{Corpus construction}

The corpus was built from two complementary types of material. First, the SIAGE
Chair provided internal action sheets collected in the context of previous work
on territorial adaptation to population ageing. Second, external institutional
and associative websites were selected because they contained practical sheets,
initiative descriptions, research-project records, or public-action resources
related to the project's scope.

Only documents containing a sufficiently substantial description of an
initiative, programme, action, or research project were retained. Navigation
pages, empty pages, duplicated pages, and pages without a meaningful connection
to territorial adaptation to ageing were excluded whenever they could be
identified automatically or during inspection.

Table~\ref{tab:final-corpus} presents the final corpus actually processed in
the project. It supersedes exploratory source lists: CNSA and Fondation Médéric
Alzheimer were considered during source exploration but are not represented as
separate source directories in the final delivered corpus.

\begin{table}[H]
\centering
\begin{tabular}{lrp{7cm}}
\toprule
Source & Documents & Nature of the material \\
\midrule
CNAV & 42 & Institutional documents and prevention-related initiatives. \\
Fondation de France & 43 & Associative or territorial initiative descriptions. \\
Fiches actions SIAGE & 60 & Internal action sheets supplied by the host institution. \\
HCFEA & 23 & Institutional reports and ageing-related public-action resources. \\
IReSP Autonomie & 131 & Research-project descriptions in the field of autonomy. \\
RFVAA & 854 & Descriptions of local age-friendly initiatives. \\
\midrule
Total & 1,153 & Heterogeneous textual corpus. \\
\bottomrule
\end{tabular}
\caption{Final corpus used in the project.}
\label{tab:final-corpus}
\end{table}

Several other websites were considered during the exploration phase, including
France Bénévolat, Fondation La Poste, Fondation Nicolas Hulot, AGURAM, and
Villages d'Avenir. They were not retained for automated collection because they
contained too few exploitable documents, insufficiently structured information,
or resources that were difficult to access and process automatically.

The corpus is not a statistically representative sample of all French
initiatives related to ageing. RFVAA contributes 854 of the 1,153 records, or
approximately 74\% of the corpus. This source imbalance must be considered when
interpreting descriptive statistics: the most frequent themes, initiative types,
or stakeholders may partly reflect the editorial coverage of RFVAA rather than
the full French landscape of territorial adaptation.

\section{Cleaning and preprocessing}

After collection, the raw text files require preprocessing before they can be
analysed. The source documents contain formatting noise arising from HTML
extraction, PDF conversion, duplicated page headers, fragmented lines,
inconsistent whitespace, and variable punctuation. The cleaning stage aims to
reduce these superficial differences while preserving the linguistic content and
the traceability of the original resource.

The cleaning operations include whitespace and line-ending normalisation, the
removal of obvious formatting artefacts, and file-name standardisation. The
resulting files are stored separately from the raw documents in the
\texttt{Cleaned\_Database} directory. This separation makes it possible to
repeat the cleaning process, compare original and transformed versions, and
avoid overwriting source material.

In addition to cleaning, a keyword-based relevance score was implemented as a
first quality-control mechanism. The score is not intended to classify documents
automatically as relevant or irrelevant. Instead, it identifies documents that
should be inspected first because their content appears unusually central or
peripheral to the project scope.

The score relies on weighted keywords associated with ageing and territorial
adaptation. Higher weights are assigned to specific concepts directly related to
the project, whereas broader contextual terms receive lower weights.

\begin{table}[H]
\centering
\begin{tabular}{lc}
\toprule
Keyword & Weight \\
\midrule
agisme & 5 \\
ville amie des aines & 5 \\
vada & 5 \\
vieillissement & 3 \\
personnes agees & 3 \\
aine & 3 \\
seniors & 3 \\
retraites & 3 \\
vieillir & 3 \\
retraite & 3 \\
age & 3 \\
intergenerationnel & 3 \\
chez soi & 1 \\
ehpad & 1 \\
mobilite & 1 \\
hebergement & 1 \\
autonomie & 1 \\
citoyen & 1 \\
fragile & 1 \\
vulnerable & 1 \\
dependant & 1 \\
isole & 1 \\
aidant & 1 \\
\bottomrule
\end{tabular}
\caption{Keyword weighting used for relevance scoring.}
\label{tab:keyword-weights}
\end{table}

For each document, the script counts the occurrences of these keywords in the
normalised text and applies a small bonus if a keyword appears in the file name.
The resulting score is normalised by the number of words in the document to make
documents of different lengths more comparable. The scripts produce ranked lists
of documents with the highest and lowest scores. Low-ranked documents can be
inspected for irrelevant or poorly extracted content; high-ranked documents
provide a quick check that the corpus includes strongly relevant resources.

This relevance score is therefore a tool for quality control and manual-review
prioritisation. It must not be interpreted as a supervised classifier, a
validated measure of relevance, or an evaluation of the factual quality of the
subsequent metadata extraction.

\section{Progressive structuring of the document database}

The project architecture uses distinct directories for each major transformation.
This modular design makes it possible to isolate a processing stage, inspect its
outputs, and identify where an error was introduced. The organisation of the
project is summarised below.

\begin{verbatim}
Index_Insight/
|-- Database/
|-- Cleaned_Database/
|-- Segmented_Database/
|-- Sectioned_Database/
|-- Entities_Database/
|-- Candidates_Database/
|-- Scored_Database/
|-- Final_Tagged_Database/
|-- LLM_Tagged_Database/
|-- Evaluation/
|-- outputs_graphs/
|-- Scripts/
|-- temp/
|-- venv/
\end{verbatim}

The \texttt{Database} directory contains raw collected text files. The
\texttt{Cleaned\_Database} contains their normalised versions. The first
experimental rule-based pipeline then produces several intermediate databases.

The \texttt{Segmented\_Database} stores paragraphs, lines, and sentences. This
representation is useful because relevant information may be expressed at
different textual levels. The \texttt{Sectioned\_Database} adds detected
document sections, for example sections corresponding to objectives, target
publics, context, stakeholders, dates, or territorial information when explicit
headings are present.

The \texttt{Entities\_Database} stores intermediate entity-extraction results.
The \texttt{Candidates\_Database} contains candidate values identified for
metadata fields. The \texttt{Scored\_Database} associates these candidates with
heuristic scores and confidence values. Finally, the
\texttt{Final\_Tagged\_Database} consolidates the selected values into a
structured rule-based JSON representation and produces a global JSON file.

This progressive pipeline has two main advantages. First, it provides
transparency: intermediate candidate values and scores can be inspected rather
than treating extraction as a single opaque operation. Second, it supports
debugging and iterative development. Its limitation is that manually designed
rules do not generalise easily across the full diversity of documents.

\section{LLM-assisted structured extraction}

To address the limitations of exclusively rule-based extraction, the project
progressively adopted prompt-based structured information extraction. The
\texttt{LLM\_Tagged\_Database} directory stores JSON records generated by a
local language model. Each record corresponds to one source document and
contains structured metadata extracted from the text.

The current RFVAA extraction script uses the local model
\texttt{qwen2.5:1.5b} through the Ollama chat API. Local inference was selected
to avoid dependence on an external commercial API, maintain control over the
processing workflow, reduce operational costs, and permit iterative prompt
experimentation. If another model generated all or part of the final database,
the model identifier and the scope of that use must be reported explicitly in
the final version of the thesis.

For every document, the extraction script dynamically constructs a French prompt
that specifies the expected JSON keys, the constraints of each field, and
controlled vocabularies. The model is instructed to produce only valid JSON,
without Markdown or explanatory text, not to invent information, and to use an
empty string or an empty list when the document does not support a reliable
value.

The expected schema contains the following fields:

\begin{itemize}
    \item \texttt{titre};
    \item \texttt{territoire};
    \item \texttt{echelle};
    \item \texttt{public\_vise};
    \item \texttt{nature\_initiative};
    \item \texttt{description\_generale};
    \item \texttt{contexte};
    \item \texttt{problematique};
    \item \texttt{solution\_envisagee};
    \item \texttt{objectifs};
    \item \texttt{porteur\_initiative};
    \item \texttt{date};
    \item \texttt{thematique};
    \item \texttt{source\_site}.
\end{itemize}

These fields constitute the initial extraction schema. The original LLM values
for \texttt{thematique} and \texttt{nature\_initiative} are subsequently
preserved and complemented by normalised fields intended for filtering and
descriptive analysis.

Several fields rely on controlled vocabularies to reduce avoidable variation.
For example, \texttt{public\_vise} is constrained to values such as
\textit{aînés}, \textit{aînés et autres}, \textit{professionnels},
\textit{autre}, or \textit{public non connu}. The field \texttt{echelle} is a
list restricted to \textit{quartier}, \textit{communale},
\textit{intercommunale}, \textit{départementale}, \textit{régionale},
\textit{nationale}, and \textit{internationale}. The model is also asked to use
short and homogeneous labels for themes and initiative types.

The current configuration uses a temperature of 0, a context window of 8,192
tokens, a maximum prediction length of 700 tokens, and an input truncation
threshold of 2,500 characters in the RFVAA extraction script. The deterministic
temperature helps make successive executions more stable. However, the
truncation threshold is a limitation because information located late in a long
document may not be sent to the model.

The extraction pipeline can be summarised as follows:

\begin{center}
Raw source $\rightarrow$ text extraction $\rightarrow$ cleaning and inspection
$\rightarrow$ constrained prompt $\rightarrow$ local LLM $\rightarrow$
structured JSON $\rightarrow$ statistics and manual review.
\end{center}

Raw model responses are archived when available. This provides an audit trail:
an unexpected value can be compared with both the prompt and the model answer.
The JSON records are also checked before use, and records that cannot be parsed
or do not satisfy the expected structure can be flagged for reprocessing or
manual inspection.

The LLM approach improves flexibility when document structures vary, but it
does not eliminate the need for validation. Outputs can contain ambiguous
interpretations, unsupported inferences, inconsistent categories, or formatting
errors. The quality of an output depends on source quality, prompt design, model
version, and the part of the document made available to the model.

\section{Post-extraction taxonomy normalisation}

Although the extraction prompt encourages short and homogeneous labels, the
fields \texttt{thematique} and \texttt{nature\_initiative} may still contain
lexical variants that refer to closely related concepts. For example, labels
such as \textit{lien social}, \textit{isolement}, and \textit{inclusion sociale}
cannot be compared reliably if they are treated as unrelated categories. A
deterministic post-processing step was therefore implemented to complement the
LLM output with controlled, documented categories.

The original fields \texttt{thematique} and \texttt{nature\_initiative} remain
unchanged. The normalisation script adds the fields
\texttt{thematique\_normalisee},
\texttt{nature\_initiative\_normalisee}, and
\texttt{version\_taxonomie}. This separation preserves the original wording
generated from the source document while providing consistent labels for
filtering, statistics, and visualisation.

The thematic taxonomy contains fifteen substantive categories:
\textit{autonomie et avancée en âge}, \textit{santé et prévention},
\textit{aidants, handicap et vulnérabilités}, \textit{lien social et lutte
contre l'isolement}, \textit{intergénérationnel}, \textit{mobilité et
accessibilité}, \textit{habitat et logement}, \textit{numérique et accès à
l'information}, \textit{participation citoyenne et gouvernance},
\textit{culture et loisirs}, \textit{activité physique et sport},
\textit{espaces publics et aménagement}, \textit{services et
accompagnement}, \textit{politiques publiques et recherche}, and
\textit{environnement et développement durable}. A record may receive several
normalised themes when its original labels explicitly refer to several
dimensions.

The initiative-type taxonomy includes \textit{atelier}, \textit{formation},
\textit{sensibilisation}, \textit{enquête ou étude}, \textit{rapport ou
publication}, \textit{plaidoyer}, \textit{outil numérique},
\textit{événement}, \textit{accompagnement}, \textit{dispositif ou
programme}, \textit{offre de service}, \textit{groupe de parole},
\textit{consultation citoyenne}, \textit{lieu ou espace dédié}, and
\textit{projet d'aménagement}. Unlike themes, each record receives one
normalised initiative type.

The normalisation is rule-based and does not call the LLM again. It operates on
the extracted labels rather than on the full source text, making the
transformation reproducible and inspectable. Values that cannot be mapped
unambiguously are assigned to \textit{autre / à revoir} and listed in a review
report. This conservative design avoids masking uncertainty behind an overly
broad automatic category and allows the taxonomy to be refined through
documented manual review.

At the time of analysis, the taxonomy was applied to 1,143 available LLM JSON
records. Ten corpus documents (five IReSP records and five RFVAA records) did
not yet have an LLM output and were consequently retained in the corpus but
excluded from extraction-based descriptive outputs.

\section{Indexing, exploration, and analytical outputs}

One objective of the project is to make the corpus easier to explore and compare.
Without structured metadata, the corpus would remain a collection of independent
text files that are difficult to navigate manually. The normalised metadata make
it possible to access consistent thematic and initiative-type categories
directly, rather than parsing the full text for every query or treating minor
lexical variations as distinct concepts.

The structured corpus can support several forms of documentary exploration:

\begin{itemize}
    \item keyword-based search in titles and descriptions;
    \item metadata filtering by source, territorial scale, target audience, or
    initiative type;
    \item thematic exploration;
    \item territorial comparison;
    \item comparison between SIAGE action sheets and other documents;
    \item navigation through stakeholders, dates, and initiative categories.
\end{itemize}

The project also contains scripts that create spreadsheets and charts from the
LLM-tagged records. These outputs include distributions of themes, initiative
types, stakeholders, territories, temporal mentions, and comparisons between
action sheets and other documents. Cost-related information is also extracted
exploratorily when it can be identified in the records.

At the current stage, the indexing system remains a structured prototype rather
than a fully operational search engine. The project produces machine-readable
JSON data and analytical outputs, but does not yet deploy a dedicated retrieval
engine, vector database, or public web interface. Consequently, the thesis
should describe semantic retrieval, embeddings, FAISS, Elasticsearch, and
interactive interfaces as future possibilities rather than as implemented
features.

Future work could include semantic search based on embeddings, hybrid retrieval
combining lexical and semantic methods, clustering of initiatives according to
thematic proximity, recommendation of similar initiatives, and an interactive
interface designed with the future users of the INSIGHT project.

\section{Analytical questions and interpretation framework}

Beyond the technical objective of structuring the corpus, the project supports
an exploratory analysis of territorial responses to population ageing. The
following questions guide this analysis:

\begin{enumerate}
    \item What are the principal types of initiatives described in the corpus?
    \item How do these types vary according to thematic domain, territorial
    scale, target audience, and initiative holder?
    \item What roles are attributed to workshops in the documents, compared with
    other forms of response such as services, digital tools, training,
    awareness-raising activities, or territorial planning projects?
    \item To what extent are dates and cost information available, structured,
    and comparable across sources?
\end{enumerate}

These questions are exploratory. The corpus documents initiatives that were
collected from selected sources; it does not provide a representative survey of
all initiatives in France, nor does it measure their effectiveness. The analysis
therefore describes how initiatives are represented in the corpus and does not
establish causal claims about their impact.

\subsection{Typology of territorial responses}

The field \texttt{nature\_initiative} makes it possible to distinguish several
forms of response to ageing-related challenges, including workshops, training,
awareness-raising activities, services, support schemes, digital tools, events,
publications, discussion groups, dedicated places, and planning projects. This
typology can be analysed jointly with \texttt{thematique},
\texttt{public\_vise}, \texttt{echelle}, and \texttt{porteur\_initiative}.
For example, it can indicate whether certain response types are more often
associated with social connection, health prevention, housing, mobility, or
digital inclusion, and whether they are mainly carried by public institutions,
associations, or local authorities.

The objective is not to rank these response types. It is to make visible the
diversity of intervention formats and to identify the combinations of problems,
publics, actors, and territorial scales that occur in the collected material.

\subsection{Interpreting the prominence of workshops}

If workshops are frequent in the corpus, this can be investigated by comparing
their thematic distribution, target publics, objectives, and stated
problematiques with those of other initiative types. The analysis can identify
the functions attributed to workshops in the documents: for instance,
prevention, information, acquisition of practical knowledge, social connection,
peer exchange, participation, or the identification of local needs.

However, frequency does not demonstrate effectiveness. The corpus generally
describes the intended function or expected benefit of a workshop; it does not
systematically contain outcome measures or control groups. The thesis should
therefore formulate conclusions as follows: workshops are represented as a
particular type of response and are associated with particular intended
mechanisms, rather than claiming that they produce a proven effect.

\subsection{Temporal and cost information}

The date field contains temporal expressions mentioned in the documents. A
frequent year, such as 2019, may refer to the start of an initiative, a
publication date, a funding period, an event date, or a date cited in the
description. It should not automatically be interpreted as the year when an
initiative began. A temporal analysis is meaningful only after distinguishing,
where possible, the semantic role of each date and reporting the proportion of
records for which this role remains uncertain.

Cost information can provide an additional exploratory perspective. The project
contains a dedicated extraction procedure for cost-related mentions, but these
values may refer to a total budget, a grant, a cost per participant, a duration,
or an amount not directly associated with the initiative. Before comparing
amounts, the analysis should report the coverage rate, currency, unit, year,
and semantic type of each value. Costs should therefore be presented as
documented information available for a subset of the corpus, not as a complete
or directly comparable economic evaluation.

\subsection{NLP and interdisciplinary contribution}

The contribution of natural language processing to this project is not to
replace social-science interpretation. Its role is to make a large, fragmented,
and heterogeneous documentary corpus tractable: it supports systematic
collection, structured description, filtering, comparison, and the
identification of patterns that would be difficult to observe by reading more
than one thousand documents individually.

The interdisciplinary value of the project arises from this complementarity.
Research on ageing, territorial policies, participation, and social inclusion
provides the concepts needed to interpret the data. NLP provides the technical
means to organise and explore the material at scale. Interpretation, category
validation, and the assessment of policy relevance remain human and
context-dependent tasks.

\section{Methodological limits}

Several limits must be considered when interpreting the corpus and the outputs.
First, the corpus is unevenly distributed across sources, with RFVAA strongly
represented. Second, ten documents did not yet have an LLM JSON output at the
time of analysis and are therefore absent from extraction-based statistics.
Third, raw web and PDF extraction may preserve boilerplate text, incomplete
metadata, duplicate passages, or encoding artefacts. Fourth,
territorial references can refer to a location, an organisation, or a project
area, while the initiative holder can be confused with a partner or funder.

Controlled vocabularies reduce output variability but may hide
meaningful distinctions between initiatives. Finally, a populated metadata field
is not necessarily factually correct. Automatic completion and rule-based
confidence indicators are useful for quality control, but they do not replace a
human reference evaluation. The next chapter therefore distinguishes technical
coverage from factual extraction accuracy and presents the evaluation strategy.
